\subsection{Statistische Grundlagen}
\shortterm \bi{Stichprobe} Die Gesamtheit der Beobachtungen $x_1, \ldots, x_n$ oder Z.V. $\cX_1, \ldots, \cX_n$. $n$ ist \bi{Stichprobenumfang}
% TODO: Think of more concise def

\shortdefinition[Parameterraum] Ein Datensatz $x_1, \ldots, x_n$ aus einer Stichprobe $\cX_1, \ldots, \cX_n$ kann durch (evtl. hochdim.)
\bi{Param.} $\vartheta$ aus \bi{Param.-Raum} $\Theta$ dargestellt werden. $\vartheta$ ist dabei oft die Param der Verteilung (bspw. $\vartheta = (\mu, \sigma^2)$ bei Normalv.).
Wir betrachten jeweils \textit{Familie von W-Räumen}, $(\Omega, \cF, \P_\vartheta)_{\vartheta \in \Theta}$.\\
{\scriptsize Wir haben normalen Grundraum plus W-Mass $\P_\vartheta$}
